\documentclass[12pt,a4paper]{article}

\usepackage[utf8]{inputenc}
\usepackage[T1]{fontenc}
\usepackage{graphicx}
\usepackage{amsmath}
\usepackage{amssymb}
\usepackage{geometry}
\usepackage{fancyhdr}
\usepackage{enumitem}
\usepackage{parskip}
\usepackage{xcolor}
\usepackage{anyfontsize}
\usepackage{subcaption}
\usepackage{booktabs}
\usepackage{array}
\usepackage{setspace}
\onehalfspacing
\usepackage[colorlinks=true]{hyperref}
\usepackage{xurl}
\usepackage{float}
\usepackage{longtable}
\usepackage{titlesec}
\usepackage{tabularx}
\usepackage{ragged2e}
\usepackage[most]{tcolorbox}
\usepackage{makecell}
\usepackage{listings}
\usepackage{caption}
\usepackage{pdflscape}
\usepackage{placeins}

% =====================================
% Colors and Styles
% =====================================
\definecolor{itiBlue}{RGB}{0,82,147}
\definecolor{absblack}{RGB}{0,0,0}
\definecolor{lightgray}{RGB}{245,245,245}
\definecolor{softblue}{RGB}{239,246,252}
\definecolor{softgreen}{RGB}{239,252,245}
\definecolor{darkgreen}{RGB}{33,110,57}
\definecolor{softorange}{RGB}{255,248,235}
\definecolor{darkorange}{RGB}{173,94,0}
\hypersetup{
    linkcolor = absblack,
    urlcolor = itiBlue,
    citecolor = itiBlue,
    colorlinks = true
}

% =====================================
% Page Layout
% =====================================
\geometry{a4paper, margin=0.78in}
\setlength{\parindent}{0pt}
\setlength{\parskip}{3pt}
\setlist[itemize]{leftmargin=0.75cm,itemsep=0pt,topsep=1pt}
\setlist[enumerate]{leftmargin=0.75cm,itemsep=0pt,topsep=1pt}
\setlength{\headheight}{15pt}
\setlength{\tabcolsep}{3.4pt}
\renewcommand{\arraystretch}{1.12}
\emergencystretch=3em
\sloppy

\pagestyle{fancy}
\fancyhf{}
\rhead{Telecom Knobs Validation}
\lhead{DIGIS GP Team}
\renewcommand{\headrulewidth}{0.4pt}
\fancyfoot[C]{%
    \vspace{4pt}%
    \rule{\textwidth}{0.4pt}\\%
    \vspace{1pt}%
    \thepage%
}

\titleformat{\section}{\color{itiBlue}\normalfont\large\bfseries}{\thesection}{0.6em}{}
\titleformat{\subsection}{\color{absblack}\normalfont\normalsize\bfseries}{\thesubsection}{0.6em}{}
\titleformat{\subsubsection}{\color{absblack}\normalfont\small\bfseries}{\thesubsubsection}{0.6em}{}

\captionsetup[table]{position=bottom,skip=4pt,font=scriptsize,labelfont=bf}
\captionsetup[figure]{position=bottom,skip=4pt,font=scriptsize,labelfont=bf}

\newcolumntype{Y}{>{\RaggedRight\arraybackslash}X}
\newcolumntype{L}[1]{>{\RaggedRight\arraybackslash}m{#1}}
\newcolumntype{P}[1]{>{\RaggedRight\arraybackslash}p{#1}}
\newcolumntype{C}[1]{>{\Centering\arraybackslash}p{#1}}

\tcbset{
    colback=softblue,
    colframe=itiBlue,
    boxrule=0.6pt,
    arc=2pt,
    left=5pt,
    right=5pt,
    top=5pt,
    bottom=5pt
}

\lstdefinestyle{cmdstyle}{
    basicstyle=\ttfamily\scriptsize,
    breaklines=true,
    columns=fullflexible,
    frame=single,
    rulecolor=\color{lightgray},
    backgroundcolor=\color{lightgray},
    xleftmargin=2pt,
    xrightmargin=2pt,
    aboveskip=5pt,
    belowskip=5pt,
    showstringspaces=false
}
\lstset{style=cmdstyle}

\newtcolorbox{infobox}[1]{
    title=#1,
    colback=softblue,
    colframe=itiBlue,
    fonttitle=\bfseries,
    boxrule=0.6pt,
    arc=2pt
}

\newtcolorbox{successbox}[1]{
    title=#1,
    colback=softgreen,
    colframe=darkgreen,
    fonttitle=\bfseries,
    boxrule=0.6pt,
    arc=2pt
}

\newtcolorbox{warningbox}[1]{
    title=#1,
    colback=softorange,
    colframe=darkorange,
    fonttitle=\bfseries,
    boxrule=0.6pt,
    arc=2pt
}

\newcommand{\fname}[1]{\texttt{#1}}
\newcommand{\colname}[1]{\texttt{#1}}
\newcommand{\flowarrow}{\ensuremath{\Longrightarrow}}

\begin{document}

\begin{titlepage}
    \centering
    \begin{figure}[H]
        \centering
        \IfFileExists{../ITI_Logo.png}{%
            
\includegraphics[width=0.68\linewidth]{../ITI_Logo}%
        }{%
            \fbox{\parbox[c][3cm][c]{0.7\linewidth}{\centering\Large ITI Logo\\[0.3cm]\normalsize Place \texttt{ITI\_Logo.png} in the required logo path}}
        }
    \end{figure}

    \vspace{1.05cm}
    \textcolor{absblack}{\rule{1\textwidth}{1.5pt}} \\[0.75cm]
    {\fontsize{21}{27}\selectfont\bfseries LTE-DL Telecom Knobs Validation Report}\\[0.35cm]
    {\fontsize{15}{18}\selectfont\bfseries Thresholds, Bins, Demand Signals, and Score Calibration}\\[0.25cm]
    \textcolor{absblack}{\rule{1\textwidth}{1.5pt}} \\[1.20cm]

    \vspace*{0.25cm}
    \begin{tabular}{@{}l@{\hspace{1.3cm}}l@{}}
        \textbf{\fontsize{14}{16}\selectfont Created by:} & \fontsize{16}{18}\selectfont DIGIS GP Team \\ [0.2cm]
        & \fontsize{12}{14}\selectfont \textcolor{itiBlue}{\href{mailto:osamaeglal143@gmail.com}{osamaeglal143@gmail.com}} \\[0.2cm]
        & \fontsize{12}{14}\selectfont \textcolor{absblack}{Intake 46} \\[0.8cm]
        \textbf{\fontsize{14}{16}\selectfont Submitted to:} & \fontsize{16}{18}\selectfont Eng. Ibrahim El Shal \\[0.8cm]
        \textbf{\fontsize{14}{16}\selectfont Report Date:} & \fontsize{16}{18}\selectfont \today \\
    \end{tabular}
    \vfill
    \textcolor{absblack}{\rule{1\textwidth}{1pt}}
    \thispagestyle{empty}
\end{titlepage}

\small
\setstretch{1.00}

\section{Report Objective}
This report is limited to \textbf{telecom knob validation} for the LTE downlink throughput anomaly detector. It does not describe the full cleaning pipeline, RCA logic, or machine-learning extension. The objective is to justify the technical thresholds and assumptions used for anomaly detection based on the provided distributions and telecom reasoning.

\begin{infobox}{Validation principle}
A knob is accepted if it is supported by the observed LTE-DL throughput distribution, radio KPI separation, RB-demand behavior, or by conservative telecom logic that avoids over-flagging low-demand and weak-reference cases.
\end{infobox}

\section{Throughput Threshold Knobs}
Figure~\ref{fig:tp_dist} is the main evidence for the absolute throughput thresholds. The distribution is right-skewed: many samples are low/medium throughput, while a long tail reaches high LTE-DL throughput values.

\begin{figure}[H]
\centering
\includegraphics[width=0.96\linewidth]{Images/01_pre_anomaly_throughput_distribution.png}
\caption{LTE-DL throughput distribution before anomaly detection. The plot shows the very-low threshold at 5 Mbps, low threshold at 10 Mbps, global P10 around 4.99 Mbps, P25 around 15.43 Mbps, and median around 38.80 Mbps.}
\label{fig:tp_dist}
\end{figure}

\begin{table}[H]
\centering
\scriptsize
\begin{tabularx}{\textwidth}{P{3.3cm}P{3.2cm}Y}
\toprule
\textbf{Knob} & \textbf{Value} & \textbf{Validation logic} \\
\midrule
Very-low throughput & $TP \leq 5$ Mbps & Aligned with the observed global P10, approximately 4.99 Mbps. Used for severe low-throughput evidence, not as the only anomaly rule. \\
Low-throughput flag & $TP < 10$ Mbps & Lies between global P10 and P25. It captures clear low-QoE samples while avoiding a too-broad 15--20 Mbps hard flag. \\
Soft score boundary & $TP < 20$ Mbps & Used only for continuous severity scoring. It allows 10--20 Mbps samples to receive mild severity if other evidence exists. \\
Global P10/P25 & P10 $\approx 4.99$, P25 $\approx 15.43$ Mbps & Used as empirical lower-tail reference values. They support using 5 Mbps as very-low and 10 Mbps as the operational low-throughput threshold. \\
Median reference & Median $\approx 38.80$ Mbps & Confirms that 10 Mbps is materially below typical LTE-DL behavior in this dataset. \\
\bottomrule
\end{tabularx}
\caption{Validated throughput threshold knobs.}
\label{tab:tp_knobs}
\end{table}

\section{Radio KPI Bin Knobs}
The RSRP and SINR plots show strong throughput separation between bins, so these are reliable radio-condition knobs. RSRQ is weaker and less monotonic; therefore it is kept as supporting evidence, not as a dominant standalone trigger.

\begin{figure}[H]
\centering
\begin{subfigure}{0.49\textwidth}
    \centering
    \includegraphics[width=\linewidth]{Images/02_throughput_by_rsrp_bin.png}
    \caption{Throughput by RSRP bin.}
\end{subfigure}
\hfill
\begin{subfigure}{0.49\textwidth}
    \centering
    \includegraphics[width=\linewidth]{Images/04_throughput_by_sinr_bin.png}
    \caption{Throughput by SINR bin.}
\end{subfigure}
\caption{RSRP and SINR validation. Median and upper-quartile throughput increase as radio quality improves, supporting the selected bin ranges.}
\label{fig:rsrp_sinr}
\end{figure}

\begin{table}[H]
\centering
\scriptsize
\begin{tabularx}{\textwidth}{P{2.6cm}P{4.8cm}Y}
\toprule
\textbf{KPI} & \textbf{Bins / flags} & \textbf{Validation and usage} \\
\midrule
RSRP & Poor $< -120$ dBm; Fair $[-120,-100)$; Good $[-100,-90)$; Excellent $\geq -90$ & Strongly supported by the RSRP boxplot. Used as a main coverage knob and for radio-limited labeling. \\
SINR & Poor $<0$ dB; Fair $[0,10)$; Good $[10,15)$; Excellent $\geq15$ & Strongly supported by the SINR boxplot. Used as a main quality/interference knob and in expected-throughput grouping. \\
RSRQ & Poor $< -20$ dB; Fair $[-20,-15)$; Good $[-15,-10)$; Excellent $\geq -10$ & Included as quality/loading evidence. Because the distribution is not fully monotonic, RSRQ is not allowed to dominate alone. \\
BLER & High BLER $>10\%$ & Telecom logic: high BLER indicates retransmissions and low useful throughput. Used as supporting PHY evidence. \\
CQI & Low CQI $<7$ & Telecom logic: low CQI indicates limited channel quality. Used as supporting evidence. \\
\bottomrule
\end{tabularx}
\caption{Validated radio and PHY knobs.}
\label{tab:radio_knobs}
\end{table}

\begin{figure}[H]
\centering
\includegraphics[width=0.80\linewidth]{Images/03_throughput_by_rsrq_bin.png}
\caption{RSRQ validation. The weaker monotonic separation justifies using RSRQ as evidence, not as a primary anomaly trigger.}
\label{fig:rsrq}
\end{figure}

\section{RB Demand and CA Knobs}
RB usage is the main demand/resource-confidence knob. Figure~\ref{fig:rb_bucket} shows a clear throughput increase from low RB usage to high RB usage, supporting the decision to analyze all valid LTE-DL rows while using RB to distinguish low-demand/resource-uncertain cases from strong underperformance cases.

\begin{figure}[H]
\centering
\includegraphics[width=0.92\linewidth]{Images/05e_throughput_by_rb_usage_bucket.png}
\caption{Throughput by RB usage/demand bucket. Throughput increases strongly with RB bucket, supporting RB as the primary demand/resource-confidence signal.}
\label{fig:rb_bucket}
\end{figure}

\begin{table}[H]
\centering
\scriptsize
\begin{tabularx}{\textwidth}{P{3.3cm}P{4.4cm}Y}
\toprule
\textbf{Knob} & \textbf{Rule} & \textbf{Validation logic} \\
\midrule
Primary RB reference & Prefer \colname{lte\_agg\_rb\_dl}; fallback to \colname{lte\_rb\_count} & \colname{lte\_agg\_rb\_dl} is the most complete aggregate RB resource indicator and best reflects scheduled DL resources. \\
RB buckets & Low, below-median, medium, high using empirical RB quantiles & Supported by the RB plot: throughput rises materially across buckets. \\
Low-RB treatment & Low throughput + low RB is uncertainty/context, not automatic severe anomaly & Low RB may mean low demand, low scheduling, or resource limitation. Score is capped unless stronger evidence exists. \\
Medium/high RB gate & Strong P75, CA, and RB-efficiency anomalies require medium/high RB & Confirms that resources were actually scheduled before calling underperformance high-confidence. \\
CA active & \colname{carrier\_count} $\geq2$ or SCell evidence exists & Telecom logic: multiple carriers should increase achievable throughput if scheduling and radio are healthy. \\
CA underperformance & CA active, medium/high RB, expected $\geq40$ Mbps, actual $<15$ Mbps, ratio $\leq0.35$, no obvious poor RSRP/SINR & Prevents confusing radio-limited or low-demand cases with CA underperformance. \\
\bottomrule
\end{tabularx}
\caption{Validated RB-demand and CA knobs.}
\label{tab:rb_ca_knobs}
\end{table}

\section{Expected-Throughput and Temporal Knobs}
The expected-throughput logic uses two references. \textbf{Radio P75} estimates achievable throughput under similar radio and CA conditions. \textbf{RB-conditioned P75} estimates achievable throughput after also accounting for RB allocation. Both are needed because radio P75 can reveal scheduling/resource issues, while RB-conditioned P75 confirms true inefficiency despite allocated resources.

\begin{table}[H]
\centering
\scriptsize
\begin{tabularx}{\textwidth}{P{3.1cm}P{4.9cm}Y}
\toprule
\textbf{Knob} & \textbf{Value / rule} & \textbf{Validation logic} \\
\midrule
P75 quantile & $q=0.75$ & Represents achievable upper-normal throughput without using unstable maxima. \\
Minimum group size & $n \geq 30$ & Conservative statistical reliability rule for group references. \\
Fallback confidence & Specific levels used as high confidence; broad/weak fallback used as context only & Prevents broad P75 groups from creating strong anomalies. \\
Strict P75 gate & Expected $\geq25$ Mbps, gap $\geq15$ Mbps, actual/expected $\leq0.40$, high-confidence P75, medium/high RB & Ensures that strong underperformance requires both expected capacity and demand/resource evidence. \\
Rolling window & 11 samples; minimum 5 samples & Sample spacing is sufficiently consecutive, so sample-based rolling is acceptable; seconds are shown in context plots. \\
Rolling/sudden drop & Drop ratio $\leq0.50$, gap $\geq5$ Mbps & Detects local throughput collapse without confusing it with sustained low-window behavior. \\
Sustained window & Rolling median below low threshold with enough samples & Captures periods that are overall low throughput, not only sudden drops. \\
Event windows & Handover $\pm5$s; RACH failure $\pm10$s; CA change $\pm10$s; any event $\pm30$s & Events are timing context only. They support RCA but do not create anomalies alone. \\
\bottomrule
\end{tabularx}
\caption{Expected-throughput, temporal, and event-context knobs.}
\label{tab:p75_temporal_knobs}
\end{table}

\section{Score and Severity Knob Validation}
Figure~\ref{fig:score_qc} validates the score and method behavior after applying the telecom knobs. The method plot shows that context flags can be frequent, while final strict methods remain controlled. The score distribution is checked to ensure the thresholds Medium 35, High 60, and Critical 80 remain usable.

\begin{figure}[H]
\centering
\begin{subfigure}{0.49\textwidth}
    \centering
    \includegraphics[width=\linewidth]{Images/11_anomaly_method_percentages.png}
    \caption{Method/case flagging rate.}
\end{subfigure}
\hfill
\begin{subfigure}{0.49\textwidth}
    \centering
    \includegraphics[width=\linewidth]{Images/10b_anomaly_score_distribution_flagged_only.png}
    \caption{Flagged-score distribution.}
\end{subfigure}
\caption{Score and flag validation. Thresholds are treated as prioritization knobs, not as probability outputs.}
\label{fig:score_qc}
\end{figure}

\begin{table}[H]
\centering
\scriptsize
\begin{tabularx}{\textwidth}{P{3.2cm}P{4.5cm}Y}
\toprule
\textbf{Score knob} & \textbf{Value / rule} & \textbf{Reason} \\
\midrule
Severity thresholds & Medium 35, High 60, Critical 80 & Practical prioritization cutoffs. Histogram is used to check if Critical is too frequent. \\
Ratio severity start & Actual/expected below 0.60 & Mild gaps are ignored; severity starts when actual is clearly below expected. \\
Gap saturation & 80 Mbps & A loss near 80 Mbps is already severe; clipping avoids high expected values dominating. \\
Low-throughput score scale & Below 20 Mbps & Soft scoring boundary; 10 Mbps remains the hard low-throughput flag. \\
Drop score scale & 40 Mbps local gap & A 40 Mbps local collapse is severe enough to saturate the drop component. \\
RB-efficiency weight & Low/moderate contribution & RB already affects demand gating, RB-conditioned P75, low-RB caps, and evidence bonuses. \\
Low-RB cap & Limit low-RB uncertain cases unless stronger evidence exists & Prevents low-demand/resource-uncertain samples from becoming false High/Critical cases. \\
Score 100 cap & Reserved for catastrophic high-confidence cases & Prevents score saturation and keeps top severity meaningful. \\
\bottomrule
\end{tabularx}
\caption{Validated scoring and severity knobs.}
\label{tab:score_knobs}
\end{table}

\begin{successbox}{Final validation conclusion}
The accepted knob set is consistent with the plotted throughput, radio, and RB behavior: 10 Mbps is a reasonable hard low-throughput trigger, RSRP/SINR bins are strongly data-supported, RSRQ is correctly kept as supporting evidence, and RB usage is justified as the main demand/resource-confidence signal. Strong anomaly labels should require reliable expected throughput and medium/high RB demand whenever possible.
\end{successbox}

\end{document}
